\documentclass[12pt,a4paper]{article}

\usepackage{amsmath,amssymb}
\usepackage{graphicx}
\usepackage[table]{xcolor}
\usepackage{geometry}
\usepackage{booktabs}
\usepackage{tabularx}
\usepackage{array}
\usepackage{enumitem}
\usepackage{listings}
\usepackage{caption}
\usepackage{xurl}
\usepackage{hyperref}

\geometry{margin=2.4cm}
\setlength{\parindent}{0pt}
\setlength{\parskip}{0.55em}
\renewcommand{\arraystretch}{1.35}
\setcounter{tocdepth}{2}

% زی‌پرشین مسیر اصلی است؛ پلی‌گلاسیا امکان ساخت در محیط‌های فاقد آن را فراهم می‌کند.
\IfFileExists{xepersian.sty}{
  \usepackage{xepersian}
  \IfFontExistsTF{Amiri}
    {\settextfont[Scale=1.05]{Amiri}}
    {\settextfont[Scale=1.00]{DejaVu Sans}}
  \IfFontExistsTF{Noto Sans}
    {\setlatintextfont{Noto Sans}}
    {\setlatintextfont{DejaVu Sans}}
}{
  \usepackage{polyglossia}
  \setmainlanguage{persian}
  \setotherlanguage{english}
  \IfFontExistsTF{Amiri}
    {\newfontfamily\persianfont[Script=Arabic,Scale=1.05]{Amiri}}
    {\newfontfamily\persianfont[Script=Arabic,Scale=1.00]{DejaVu Sans}}
  \IfFontExistsTF{Noto Sans}
    {\newfontfamily\englishfont{Noto Sans}}
    {\newfontfamily\englishfont{DejaVu Sans}}
  \newcommand{\lr}[1]{\textenglish{##1}}
}

\newcommand{\eng}[1]{\lr{#1}}
\newcommand{\code}[1]{\lr{\texttt{#1}}}
\newcommand{\project}{\eng{Micro-ESPectre}}

\definecolor{codebg}{RGB}{247,247,247}
\definecolor{codegreen}{RGB}{0,110,70}
\definecolor{codepurple}{RGB}{120,40,140}

\lstset{
  basicstyle=\small\ttfamily,
  backgroundcolor=\color{codebg},
  frame=single,
  rulecolor=\color{gray!35},
  breaklines=true,
  postbreak=\mbox{\textcolor{red}{$\hookrightarrow$}\space},
  keywordstyle=\color{blue!70!black}\bfseries,
  commentstyle=\color{codegreen},
  stringstyle=\color{codepurple},
  showstringspaces=false,
  columns=fullflexible,
  keepspaces=true,
  numbers=left,
  numberstyle=\scriptsize\color{gray},
  xleftmargin=1.5em
}

\hypersetup{
  colorlinks=true,
  linkcolor=blue!55!black,
  urlcolor=cyan!45!black,
  citecolor=blue!55!black,
  pdftitle={طرح جامع توسعه دیتاست و الگوریتم‌های هوش مصنوعی بر بستر Micro-ESPectre},
  pdfauthor={امیرعلی یزدانی}
}

\title{\textbf{طرح جامع توسعه، جمع‌آوری دیتاست و الگوریتم‌های هوش مصنوعی بر بستر \project}}
\author{امیرعلی یزدانی}
\date{\today}

\begin{document}

\maketitle

\begin{abstract}
این سند، مسیر توسعه سامانه حسگری مبتنی بر اطلاعات وضعیت کانال وای‌فای را بر بستر \project ارائه می‌کند. هدف، گسترش سامانه از آشکارسازی حرکت به شمارش افراد، تشخیص ناحیه حضور، تشخیص حالت بدنی، تشخیص سقوط و پایش تنفس است. برای هر قابلیت، روش جمع‌آوری داده، الگوریتم پایه، مسیر ارتقا، معیار ارزیابی و محل اجرای نهایی مشخص شده است.
\end{abstract}

\tableofcontents
\newpage

\section{مقدمه، دامنه و وضعیت فعلی سامانه}

زیرساخت دریافت پایدار داده‌های وضعیت کانال \eng{Channel State Information (CSI)} روی \eng{ESP32-S3} در بستر \project راه‌اندازی شده است. داده‌ها از طریق \eng{MQTT} به سیستم میزبان منتقل می‌شوند و خروجی آشکارساز حرکت نیز به‌صورت بلادرنگ قابل مشاهده است. مسیر بعدی پروژه، استفاده از همین زیرساخت برای توسعه قابلیت‌های پیشرفته مبتنی بر ماتریس‌های زمانی و فرکانسی \eng{CSI} است.

معماری توسعه شامل دریافت و پردازش اولیه روی برد، ذخیره‌سازی و آموزش مدل روی سیستم میزبان و انتقال مدل‌های سبک و نهایی به لبه خواهد بود.

\subsection{ساختار داده خام}

برای بسته شماره $n$ و زیرحامل شماره $k$، پاسخ مختلط کانال به‌شکل زیر تعریف می‌شود:

\begin{equation}
H_{n,k}=I_{n,k}+jQ_{n,k}.
\end{equation}

دامنه و فاز از روابط زیر به‌دست می‌آیند:

\begin{equation}
A_{n,k}=\sqrt{I_{n,k}^{2}+Q_{n,k}^{2}},
\qquad
\phi_{n,k}=\operatorname{atan2}\!\left(Q_{n,k},I_{n,k}\right).
\end{equation}

تابع \eng{atan2} ربع صحیح فاز را حفظ می‌کند و در حالت $I=0$ نیز قابل استفاده است. در حالت \eng{HT20}، آرایش خام شامل ۶۴ مؤلفه است. پس از حذف مؤلفه‌های نامعتبر، زیرحامل‌های فعال بر اساس نگاشت واقعی فریمور انتخاب می‌شوند. همراه هر برداشت نیز زمان، شناسه گره، \eng{RSSI}، نسخه فریمور، نرخ واقعی دریافت، محیط، جلسه و برچسب کلاس ذخیره خواهد شد.

\section{فرآیند جمع‌آوری و برچسب‌گذاری دیتاست}

\subsection{مسیر دریافت و قالب ذخیره‌سازی}

اسکریپت جمع‌آوری روی سیستم میزبان به تاپیک \eng{MQTT} متصل می‌شود و داده‌ها را در قالب \eng{NPY}، \eng{NPZ} یا یک پایگاه داده ساختاریافته ذخیره می‌کند. هر رکورد شامل بردار \eng{CSI}، زمان، شماره توالی، شناسه گره، \eng{RSSI} و برچسب سناریو خواهد بود.

هر نمونه آموزشی یک تانسور پنجره‌ای است:

\begin{equation}
\mathbf{X}\in\mathbb{R}^{T\times K\times C},
\end{equation}

که در آن $T$ تعداد گام‌های زمانی، $K$ تعداد زیرحامل‌های معتبر و $C$ تعداد کانال‌های ورودی مانند دامنه و فاز است.

\subsection{نرخ نمونه‌برداری و پنجره‌ها}

تنظیمات پایه جمع‌آوری داده به‌صورت زیر در نظر گرفته می‌شود:

\begin{itemize}
  \item نرخ نمونه‌برداری ۵۰ فریم بر ثانیه؛
  \item پنجره ۲۰۰ بسته‌ای، معادل چهار ثانیه؛
  \item حداقل ۱۰٬۰۰۰ فریم برای هر کلاس و سناریو؛
  \item پنجره ۳۰ تا ۶۰ ثانیه برای تنفس؛
  \item هم‌پوشانی ۵۰ درصد برای تولید نمونه‌های آموزشی.
\end{itemize}

برای افزایش قابلیت تعمیم، این داده‌ها در چند جلسه، زمان، چیدمان و ترکیب انسانی متفاوت جمع‌آوری می‌شوند.

\subsection{تقسیم آموزش، اعتبارسنجی و آزمون}

داده‌ها بر اساس جلسه و سوژه به سه بخش آموزش، اعتبارسنجی و آزمون نهایی تقسیم می‌شوند. تمام پنجره‌های متعلق به یک جلسه فقط در یکی از این بخش‌ها قرار می‌گیرند. آزمون‌های \eng{Leave-One-Subject-Out} و \eng{Leave-One-Room-Out} نیز برای ارزیابی عملکرد روی افراد و محیط‌های جدید اجرا خواهند شد.

\section{کنترل کیفیت و پیش‌پردازش}

پیش از آموزش، بسته‌های تکراری یا خارج از ترتیب شناسایی می‌شوند و نرخ دریافت، درصد بسته‌های مفقود، طول بردار \eng{CSI} و هم‌ترازی زمانی برچسب‌ها کنترل می‌گردد. دامنه در سطح پنجره نرمال‌سازی می‌شود و برای فاز نیز بازکردن فاز، حذف روند و کاهش جهش‌های ناگهانی انجام خواهد شد.

\subsection{ویژگی‌های پایه}

ویژگی‌های پایه برای مدل‌های کلاسیک عبارت‌اند از:

\begin{enumerate}
  \item میانگین، انحراف معیار، دامنه میان‌چارکی، چولگی و کشیدگی؛
  \item هم‌بستگی زمانی، انرژی باندی و آنتروپی طیفی؛
  \item مقادیر ویژه یا تکین ماتریس پنجره به‌عنوان توصیف‌کننده ساختار غالب؛
  \item تغییر شکل کانال، اختلاف زمانی و ویژگی‌های مقاوم به تغییر مقیاس.
\end{enumerate}

مجموعه نهایی ویژگی‌ها با استفاده از داده اعتبارسنجی و بر اساس عملکرد هر قابلیت انتخاب خواهد شد.

\section{شمارش افراد در محیط}

\subsection{صورت‌بندی و خط پایه}

هدف، تخمین تعداد افراد از صفر تا ده نفر است. دو خط پایه به‌صورت موازی ارزیابی می‌شوند:

\begin{itemize}
  \item دسته‌بندی چندکلاسه با \eng{SVM} و \eng{Random Forest}؛
  \item رگرسیون یا رگرسیون ترتیبی با خروجی محدودشده به بازه صفر تا ده.
\end{itemize}

در این مسئله، ساختار ترتیبی برچسب‌ها مهم است؛ اشتباه بین چهار و پنج نفر از اشتباه بین صفر و ده نفر کم‌هزینه‌تر است. معیارهای اصلی شامل خطای مطلق میانگین \eng{MAE}، دقت دقیق، دقت با تحمل یک نفر، ماتریس درهم‌ریختگی و عملکرد به تفکیک محیط هستند.

\subsection{پروتکل جمع‌آوری}

برای هر تعداد، افراد به‌صورت گام‌به‌گام وارد محیط می‌شوند و داده کلاس‌های صفر تا ده نفر ثبت می‌گردد. این آزمایش در چند جلسه و با تغییر ترکیب افراد، محل ایستادن یا نشستن و سطح حرکت تکرار خواهد شد. حالت‌های سکون، حرکت کم و حرکت معمول نیز به‌صورت جداگانه برچسب‌گذاری می‌شوند.

\subsection{مدل ارتقایافته}

پس از تثبیت خط پایه، یک شبکه کوچک \eng{1D-CNN} روی دنباله زمانی یا یک \eng{2D-CNN} روی ماتریس زمان و فرکانس بررسی می‌شود. مدل نهایی بر اساس عملکرد آن روی افراد و محیط‌های دیده‌نشده انتخاب خواهد شد.

\section{تشخیص ناحیه و مکان‌یابی فضایی}

\subsection{خط پایه اثرانگشت‌نگاری}

محیط به چند ناحیه کاربردی تقسیم می‌شود و برای هر ناحیه، اثرانگشت \eng{CSI} در جلسه‌های مستقل ثبت می‌گردد. طبقه‌بند \eng{k-NN} با فاصله اقلیدسی، کوسینوسی یا ماهالانوبیس خط پایه مناسبی است. با یک لینک منفرد، تشخیص چند ناحیه در یک محیط ثابت ممکن است، اما به چیدمان و محیط حساس است و نباید مکان‌یابی هندسی عمومی تلقی شود.

\subsection{ارتقا با چند لینک}

افزودن گره‌ها یا لینک‌های مستقل، تنوع دید رادیویی را افزایش می‌دهد و ابهام برخی نواحی را کاهش می‌دهد، اما آن را به‌طور کامل حذف نمی‌کند. برای ادغام چند لینک، زمان‌ها باید هم‌تراز شوند و کیفیت هر لینک به‌صورت جداگانه ثبت گردد. مدل باید توان ادامه کار در صورت قطع یک گره را نیز ارزیابی کند.

\subsection{محدودیت تخمین زاویه ورود}

الگوریتم‌هایی مانند \eng{MUSIC} به آرایه چندآنتنه همدوس، هندسه معلوم، کالیبراسیون فاز و معمولاً چند جریان فضایی نیاز دارند. چند برد مستقل \eng{ESP32-S3} بدون هم‌زمان‌سازی فازی جایگزین مستقیم چنین آرایه‌ای نیستند. بنابراین \eng{AoA/MUSIC} فقط در صورت استفاده از سخت‌افزار مناسب مانند کارت شبکه پژوهشی چندآنتنه یا آرایه همدوس، به‌عنوان مسیر جداگانه بررسی می‌شود.

\section{تشخیص حالت بدنی و سقوط}

\subsection{تفکیک دو مسئله}

تشخیص حالت ایستاده، نشسته و درازکش یک مسئله طبقه‌بندی حالت است؛ تشخیص سقوط یک مسئله رویداد زمانی است. قرار دادن همه آن‌ها در یک طبقه‌بند تک‌فریمی می‌تواند نرخ مثبت کاذب بالایی ایجاد کند. خط پایه پیشنهادی دو مرحله دارد:

\begin{enumerate}
  \item تشخیص حرکت سریع یا تغییر شدید در پنجره کوتاه؛
  \item تأیید حالت پس از رویداد و دوره بی‌حرکتی در پنجره بلندتر.
\end{enumerate}

مدل پایه می‌تواند \eng{Random Forest} یا \eng{MLP} روی ویژگی‌های یک پنجره زمانی باشد. مدل ارتقایافته \eng{CNN-LSTM} یا شبکه کانولوشن زمانی است. طیف‌نگاره داپلر حاصل از \eng{STFT} نیز ورودی مکمل مفیدی است.

\subsection{طراحی داده و ارزیابی ایمن}

سقوط واقعی نباید بدون پروتکل ایمنی اجرا شود. سقوط‌های شبیه‌سازی‌شده، نشستن سریع، درازکشیدن عادی، برداشتن جسم از زمین و افتادن اشیا باید در داده حضور داشته باشند. ارزیابی باید رویدادمحور باشد و حساسیت، دقت مثبت، تعداد هشدار کاذب در ساعت و تأخیر تشخیص را گزارش کند. این سامانه پژوهشی جایگزین تجهیز پزشکی یا سامانه اضطراری تأییدشده نیست.

\section{پایش تنفس و علائم حیاتی}

حرکت قفسه سینه در تنفس می‌تواند تغییرات کم‌دامنه‌ای در مسیرهای چندگانه ایجاد کند. بازه متداول بررسی تنفس بزرگسالان حدود ۰٫۱ تا ۰٫۵ هرتز است، اما بازه نهایی باید متناسب با جمعیت هدف و مرجع اندازه‌گیری تعیین شود.

\subsection{خط پایه پردازش سیگنال}

\begin{enumerate}
  \item بازنمونه‌برداری روی شبکه زمانی یکنواخت و رد پنجره‌های کم‌کیفیت؛
  \item انتخاب یا ادغام زیرحامل‌ها بر اساس پایداری و انرژی باند تنفس؛
  \item فیلتر میان‌گذر با کنترل اثر مرزی؛
  \item تخمین نرخ با قله‌یابی، عبور از صفر یا بیشینه طیفی؛
  \item هموارسازی زمانی همراه با گزارش شاخص اطمینان.
\end{enumerate}

انتخاب «بهترین زیرحامل» باید فقط با داده آموزش یا بخش آغازین همان جلسه انجام شود تا نشت اطلاعات از آزمون رخ ندهد. فاز خام همیشه از دامنه بهتر نیست و باید پس از پالایش و در مقایسه مستقیم ارزیابی شود.

\subsection{مرجع واقعی و معیارها}

برای برچسب معتبر، استفاده از کمربند تنفسی یا مرجع هم‌زمان توصیه می‌شود. شمارش دستی برای آزمایش مقدماتی قابل استفاده است، اما برای ارزیابی دقیق کافی نیست. معیارها شامل خطای مطلق نرخ تنفس، پوشش خروجی معتبر، تأخیر، عملکرد بر حسب فاصله و نرخ شکست در حضور حرکت هستند.

استخراج ضربان قلب در بازه تقریبی ۰٫۸ تا ۲ هرتز از \eng{CSI} تک‌لینک و سخت‌افزار ارزان بسیار دشوارتر از تنفس است. این قابلیت تا زمان مقایسه با مرجع هم‌زمان و آزمون مستقل، باید هدف پژوهشی پرریسک تلقی شود. روش‌هایی مانند \eng{VMD} یا \eng{EMD} ممکن است مفید باشند، اما جایگزین کنترل کیفیت و مرجع واقعی نیستند.

\section{طراحی آزمایش و معیارهای پذیرش}

\begin{table}[htbp]
\centering
\small
\caption{حداقل طراحی ارزیابی برای قابلیت‌های هدف}
\label{tab:evaluation}
\begin{tabularx}{\textwidth}{>{\raggedleft\arraybackslash}p{2.3cm} X X X}
\toprule
\textbf{قابلیت} & \textbf{تفکیک ضروری داده} & \textbf{معیارهای اصلی} & \textbf{آزمون تعمیم} \\
\midrule
شمارش افراد & جلسه، ترکیب افراد و محیط & \eng{MAE}، دقت دقیق و تحمل یک نفر & سوژه و اتاق دیده‌نشده \\
تشخیص ناحیه & جلسه، چیدمان و ناحیه & دقت متوازن و ماتریس درهم‌ریختگی & روز و چیدمان دیده‌نشده \\
حالت بدنی & سوژه، حالت و زاویه نسبت به لینک & \eng{Macro-F1} و فراخوانی هر کلاس & سوژه دیده‌نشده \\
سقوط & رویداد، فعالیت مشابه و عدم رویداد & حساسیت، دقت مثبت، هشدار کاذب در ساعت و تأخیر & سوژه و سناریوی دیده‌نشده \\
تنفس & سوژه، فاصله، حالت و کیفیت سیگنال & خطای نرخ، پوشش معتبر و نرخ شکست & سوژه و موقعیت دیده‌نشده \\
\bottomrule
\end{tabularx}
\end{table}

هر نتیجه باید همراه با فاصله اطمینان و تعداد جلسه‌های مستقل گزارش شود. گزارش صرفِ دقت پنجره‌ای روی پنجره‌های هم‌پوشان یک جلسه می‌تواند عملکرد را به‌شدت بیش‌برآورد کند.

\section{مقایسه مدل‌ها و محل اجرا}

\begin{table}[htbp]
\centering
\small
\caption{مقایسه فنی قابلیت‌های قابل توسعه بر پایه \project}
\label{tab:ai_features}
\begin{tabularx}{\textwidth}{>{\raggedleft\arraybackslash}p{2.15cm} >{\raggedleft\arraybackslash}p{2.45cm} >{\raggedleft\arraybackslash}p{2.55cm} >{\raggedleft\arraybackslash}p{2.0cm} X}
\toprule
\textbf{قابلیت هدف} & \textbf{مدل پایه} & \textbf{مدل ارتقایافته} & \textbf{محل اجرای پیشنهادی} & \textbf{نوع ورودی} \\
\midrule
شمارش افراد & \eng{SVM / Random Forest} & رگرسیون ترتیبی یا \eng{CNN} کوچک & میزبان؛ سپس مدل فشرده & دامنه مقاوم به بهره و ویژگی زمانی \\
تشخیص ناحیه & \eng{k-NN Fingerprinting} & ادغام چند لینک و شبکه کوچک & میزبان یا دروازه & اثرانگشت چندزیرحاملی و کیفیت لینک \\
حالت و سقوط & مدل پنجره‌ای کلاسیک & \eng{TCN / CNN-LSTM} & میزبان & دامنه، فاز پالایش‌شده و طیف‌نگاره \\
پایش تنفس & میان‌گذر و تخمین قله یا طیف & ادغام زیرحامل و \eng{VMD} & میزبان؛ خط پایه سبک روی برد & دامنه یا فاز پالایش‌شده همراه شاخص کیفیت \\
\bottomrule
\end{tabularx}
\end{table}

\section{نقشه راه استقرار لبه‌ای}

\begin{enumerate}
  \item \textbf{تثبیت قرارداد داده:} نسخه‌بندی طرح‌واره، زمان، فراداده و نگاشت زیرحامل‌ها؛
  \item \textbf{ساخت خط پایه میزبان:} آموزش با \eng{Scikit-learn} و پیاده‌سازی خط لوله قابل بازپخش؛
  \item \textbf{اعتبارسنجی مستقل:} آزمون گروهی بر حسب جلسه، سوژه و محیط، سپس آزمون نگه‌داشته‌شده؛
  \item \textbf{پروفایل منابع:} اندازه مدل، حافظه پایدار، حافظه موقت، زمان پردازش هر بسته و زمان استنتاج؛
  \item \textbf{فشرده‌سازی:} کوانتیزه‌سازی \eng{INT8}، هرس یا تقطیر فقط در صورت حفظ معیارهای پذیرش؛
  \item \textbf{انتخاب زمان‌اجرای مناسب:} استفاده از \eng{TensorFlow Lite Micro}، \eng{ESP-DL} یا پیاده‌سازی مستقیم \eng{C++} بر اساس پشتیبانی واقعی عملیات مدل؛
  \item \textbf{آزمون روی برد:} بررسی پایداری چندساعته، حذف بسته، دما، بازراه‌اندازی و تداخل شبکه؛
  \item \textbf{حالت ایمن:} اگر کیفیت ورودی پایین است، سامانه باید «نامعتبر یا نامطمئن» گزارش کند و خروجی قطعی تولید نکند.
\end{enumerate}

مدل‌های سنگین مانند \eng{CNN-LSTM}، \eng{MUSIC} و تجزیه‌های تکراری ابتدا روی میزبان اجرا می‌شوند. پس از اندازه‌گیری حافظه و زمان اجرا، مدل‌های مناسب کوانتیزه و به \eng{ESP32-S3} منتقل خواهند شد. برای مدل‌های آموزش‌دیده در \eng{PyTorch} نیز مسیر تبدیل، سازگاری عملگرها و دقت پس از کوانتیزه‌سازی بررسی می‌شود.

\section{جمع‌بندی و اولویت اجرایی}

اولویت نخست، ساخت دیتاست قابل اعتماد و خط پایه قابل بازپخش است. پیشنهاد می‌شود توسعه با یک آزمایش کوچک برای سه وظیفه آغاز شود: تشخیص حرکت به‌عنوان کنترل مثبت، تشخیص چهار ناحیه و شمارش صفر تا سه نفر. پس از اثبات تعمیم میان جلسه‌ها و افراد، دامنه شمارش افزایش یابد و سپس حالت بدنی و تنفس بررسی شوند. تشخیص سقوط، ضربان قلب و تخمین زاویه ورود به دلیل نیازهای ایمنی یا سخت‌افزاری، در مراحل بعدی و با مرجع مستقل قرار می‌گیرند.

\section*{منابع فنی مبنا}
\addcontentsline{toc}{section}{منابع فنی مبنا}

\begin{latin}
\begin{enumerate}
  \item ESPectre repository and documentation: \url{https://github.com/francescopace/espectre}
  \item Micro-ESPectre runtime documentation: \url{https://github.com/francescopace/espectre/blob/main/src/python/micro_espectre/README.md}
  \item ESPectre algorithms and data collection guides: \url{https://github.com/francescopace/espectre/tree/main/docs}
  \item RuView repository and advanced Wi-Fi sensing modules: \url{https://github.com/ruvnet/ruview}
  \item Internal Phase 5 report supplied with this document.
\end{enumerate}
\end{latin}

\end{document}
